\section*{Exercises on \S\arabic{exercisegroup}}
\phantomsection
\addcontentsline{toc}{section}{Exercises on \protect\S\arabic{exercisegroup}}

\begin{enumerate}
\def\labelenumi{\arabic{enumi})}
\item
  If the series with general term \(u_{n} \geqslant 0\) converges, so does the series with general term \(\sqrt{u_n u_{n+1}}\) . The converse is false in general; show that it holds if the series \((u_n)\) is decreasing.
\item
  Let \((p_{n})\) be an increasing sequence of numbers \textgreater{} 0 tending to \(+\infty\) .
\end{enumerate}

\begin{enumerate}
\def\labelenumi{\alph{enumi})}
\tightlist
\item
  If the ratio \(p_n / p_{n-1}\) tends to 1 as \(n\) increases indefinitely show that
\end{enumerate}

\[
\sum_ {k = 1} ^ {n} \frac {p _ {k} - p _ {k - 1}}{p _ {k}} \sim \log p _ {n}
\]

\[
\sum_ {k = 1} ^ {n} p _ {k} ^ {\rho} (p _ {k} - p _ {k - 1}) \sim \frac {1}{\rho + 1} p _ {n} ^ {\rho + 1} \quad \text { for } \rho > - 1
\]

(apply prop. 2 of V, p.~237).

\begin{enumerate}
\def\labelenumi{\alph{enumi})}
\setcounter{enumi}{1}
\item
  Without the hypothesis on the ratio \(p_{n}/p_{n-1}\) , show that the series with general term \((p_{n}-p_{n-1})/p_{n}\) always has an infinite sum (distinguish the two cases according to whether \(p_{n}/p_{n-1}\) tends to 1 or not).
\item
  Show that for every \(\rho > 0\) the series with general term \((p_n - p_{n-1}) / p_n p_{n-1}^\rho\) converges (compare to the series with general term \(\frac{1}{p_{n-1}^\rho} - \frac{1}{p_n^\rho}\) ).
\end{enumerate}

\begin{enumerate}
\def\labelenumi{\arabic{enumi})}
\setcounter{enumi}{2}
\item
  Show that, for every convergent series \((u_{n})\) with terms \(>0\) , there exists a series \((v_{n})\) with infinite sum, and terms \(>0\) , such that \(\liminf v_n / u_n = 0\) .
\item
  Let \((u_{n})\) be a decreasing sequence of numbers \textgreater0; if there exists an integer k such that \(ku_{kn} \geqslant u_{n}\) for all n from some point on, then the series with general term \(u_{n}\) has an infinite sum.
\item
  Let \((u_{n})\) be a series with terms \(>0\) from some point on. Show that if
\end{enumerate}

\[
\operatorname * {l i m s u p} _ {n \to \infty} \left(\frac {u _ {n + 1}}{u _ {n}}\right) ^ {n} <   \frac {1}{e}
\]

then the series converges, and if \(\liminf_{n\to\infty}\left(\frac{u_{n+1}}{u_n}\right)^n\geqslant\frac{1}{e}\) the series has an infinite sum.

\begin{enumerate}
\def\labelenumi{\arabic{enumi})}
\setcounter{enumi}{5}
\item
  Let \((u_{n})\) be a series of reals \(\neq 0\) such that \(\lim_{n\to \infty}u_n = 0\) . Show that if there exists a number \(r\) such that \(0 < r < 1\) and if from some point on one has \(-1\leqslant u_{n + 1} / u_n\leqslant r\) , then the series \((u_{n})\) converges.
\item
  Let \((u_{n})\) be a convergent series with terms \(>0\) .
\end{enumerate}

\begin{enumerate}
\def\labelenumi{\alph{enumi})}
\tightlist
\item
  Show that
\end{enumerate}

\[
\lim _ {n \rightarrow \infty} \frac {u _ {1} + 2 u _ {2} + \cdots + n u _ {n}}{n} = 0.
\]

\begin{enumerate}
\def\labelenumi{\alph{enumi})}
\setcounter{enumi}{1}
\tightlist
\item
  Show that
\end{enumerate}

\[
\sum_ {n = 1} ^ {\infty} \frac {u _ {1} + 2 u _ {2} + \cdots + n u _ {n}}{n (n + 1)} = \sum_ {n = 1} ^ {\infty} u _ {n}
\]

\[
\left(\text { write } \frac {1}{n (n + 1)} = \frac {1}{n} - \frac {1}{n + 1}\right).
\]

\begin{enumerate}
\def\labelenumi{\alph{enumi})}
\setcounter{enumi}{2}
\tightlist
\item
  Deduce from a) and b) that
\end{enumerate}

\[
\lim _ {n \to \infty} n (u _ {1} u _ {2} \dots u _ {n}) ^ {1 / n} = 0
\]

and that

\[
\sum_ {n = 1} ^ {\infty} \frac {(n ! u _ {1} u _ {2} \dots u _ {n}) ^ {1 / n}}{n + 1} <   \sum_ {n = 1} ^ {\infty} u _ {n}
\]

(apply the inequality (8) of III, p.~93).

8) a) Let \((z_{n})\) be a sequence of complex numbers. Show that if the series with general terms \(z_{n}, z_{n}^{2}, \ldots, z_{n}^{q-1}, |z_{n}|^{q}\) converge, then the infinite product with general factor \(1 + z_{n}\) converges.

\begin{enumerate}
\def\labelenumi{\alph{enumi})}
\setcounter{enumi}{1}
\tightlist
\item
  If \(z_{n}\) is real and if the series with general term \(z_{n}\) converges, then the infinite product with general factor \(1 + z_{n}\) converges if the series with general term \(z_{n}^{2}\) converges, and one has
\end{enumerate}

\[
\underset {n = 1} {\overset {\infty} {\mathrm{P}}} (1 + z _ {n}) = 0 \quad \text { if } \quad \sum_ {n = 1} ^ {\infty} z _ {n} ^ {2} = + \infty .
\]

\begin{enumerate}
\def\labelenumi{\alph{enumi})}
\setcounter{enumi}{2}
\tightlist
\item
  For every integer \(p \geqslant 2\) put
\end{enumerate}

\[
k _ {p} = [ \log \log p ], \quad h _ {p} = \sum_ {j = 1} ^ {p - 1} k _ {j}, \quad \omega_ {p} = \frac {2 \pi}{k _ {p}}.
\]

Define the sequence \((z_{n})\) of complex numbers in the following way: for \(n = h_{p} + m\) , \(0 \leqslant m \leqslant k_{p} - 1\) , put \(z_{n} = e^{mi\omega_{p}} / \log p\) . Show that for every integer q \textgreater{} 0 the series with general term \(z_{n}^{q}\) converges, but that the product with general factor \(1 + z_{n}\) does not converge.

\begin{enumerate}
\def\labelenumi{\arabic{enumi})}
\setcounter{enumi}{8}
\tightlist
\item
  \begin{enumerate}
  \def\labelenumii{\alph{enumii})}
  \tightlist
  \item
    Show, with the help of Stirling's formula, that the maximum of the function \(f_{n}(x) = \left|e^{-2x} - e^{-x}\sum_{k=0}^{n}(-1)^{k}\frac{x^{k}}{k!}\right|\) over the interval \([0, +\infty[\) tends to 0 with \(1/n\) .
  \end{enumerate}
\end{enumerate}

\begin{enumerate}
\def\labelenumi{\alph{enumi})}
\setcounter{enumi}{1}
\tightlist
\item
  Deduce, by induction on the integer p, that for every \(\varepsilon > 0\) there exists a polynomial \(g(x)\) such that \(\left|e^{-px} - e^{-x}g(x)\right| \leqslant \varepsilon\) for every \(x \geqslant 0\) (replace x by px/2 in a), and apply the induction hypothesis to \(e^{-(p-1)x/2}\) .
\end{enumerate}

\begin{enumerate}
\def\labelenumi{\arabic{enumi})}
\setcounter{enumi}{9}
\tightlist
\item
  For every number \(\alpha > 0\) derive the formula
\end{enumerate}

\[
1 ^ {\alpha n} + 2 ^ {\alpha n} + \dots + n ^ {\alpha n} \sim \frac {n ^ {\alpha n}}{1 - e ^ {- \alpha}}
\]

as n tends to \(+\infty\) .

11) For every number \(\alpha > 0\) derive the formula

\[
(1!) ^ {- \alpha / n} + (2!) ^ {- \alpha / n} + \dots + (n!) ^ {- \alpha / n} \sim \frac {1}{\alpha} \frac {n}{\log n}
\]

as \(n\) tends to \(+\infty\) (compare each term \((p!)^{-\alpha / n}\) to \(n^{-\alpha p / n}\) ).

